\scanpage{103}
\mspage{3.5}
% archive sheet number: 102
follows from lemma 3.4 that $e(\overset{\handwritten{\sim}}{S}/\underline p\mid D)$ is an integer for all $\underline p$ with $\underline p\cap R=0$ and consequently $e(S/R)=e(\overset{\handwritten{\sim}}{S}/D)=\handwritten{\sum}_{\underline p}L(\overset{\handwritten{\sim}}{S}\underline p)e(\overset{\handwritten{\sim}}{S}/\underline p\mid D)$ is an integer.

Theorem 3.8: Let $R\handwritten{\subset}S$ be a local algebra of local domains with quotient fields $k\handwritten{\subset}K$ and assume that $S$ is integrally closed. Then $r(S/R)=\handwritten{\sum}_{\underline p}L(S^f_{\underline p})e(S^f/\underline p\mid S)$ and $e(S^f/\underline p\mid S)$ are all nonzero integers, where $\underline p$ ranges through the minimal primes in $S^f$ such that $\underline p\cap R=0$. In the above formula, $L(S^f_{\underline p})=1$ for all $\underline p$ with $\underline p\cap R=0$ if and only if $K$ is separably algebraic over $k$.

\uline{Proof}: From lemma 3.6, we have $r(S/R)=e(S^f/S)=\handwritten{\sum}_{\underline p}L(S^f_{\underline p})e(S^f/\underline p\mid S)$ where $\underline p$ ranges through the minimal prime ideals in $S^f$ such that $\dim S^f/\underline p=\dim S^f$. Put $\underline p_0=\mathrm{Ker}(S^f\longrightarrow S)$. Then $\underline p_0$ is a prime ideal in $S^f$ such that $\underline p_0\cap S=0$. Therefore we are in the position to apply lemma 3.5 since $S^f$ is a local $S$-algebra and $S$ is integrally closed. Therefore $\dim S^f/\underline p=\dim S^f$ if and only if $\underline p\cap S=0$. However, it is clear that $\underline p\cap S=0$ if and only if $\underline p\cap R=0$ since $S$ is a localization of an integral extension of $R$. Since $S$ is integrally closed and $S^f/\underline p$ and $S$ have the same residue class field, it follows from proposition 3.7 that $e(S^f/\underline p\mid S\handwritten{)}\unsure{'}$ is an integer. This proves the first assertion. Now it is clear that $S^f_{\underline p_0}=K^f_k$ and hence $L(S^f_{\underline p_0})=L(K^f_k)=[K:k]_i$ by corollary 2.9. Thus if $L(S^f_{\underline p})=1$ for all $\underline p$ with $\underline p\cap R=0$, then in particular we must have that $L(S^f_{\underline p_0})=[K:k]_i=1$, i.e. $K$ is separably algebraic over $k$. Conversely if $K$ is separably algebraic over $k$ then $K\mathbin{\ooalign{\hfil$\mathrm{x}$\hfil\cr\hfil$\handwritten{\bigcirc}$\hfil\cr}}_k K$ is a separable $k$-algebra and consequently $(K\mathbin{\ooalign{\hfil$\mathrm{x}$\hfil\cr\hfil$\handwritten{\bigcirc}$\hfil\cr}}_k K)_{\underline P}$ is a field for all prime ideal $\underline P$ in $K\mathbin{\ooalign{\hfil$\mathrm{x}$\hfil\cr\hfil$\handwritten{\bigcirc}$\hfil\cr}}_k K$. Consequently $L(S^f_{\underline p})=1$ for all $\underline p$ with $\underline p\cap R=0$.
